\documentclass[10pt,a4paper]{article}
\usepackage[utf8]{inputenc}
\usepackage[T1]{fontenc}
\usepackage[spanish,es-tabla]{babel}
\usepackage{mathptmx}
\usepackage[margin=1.8cm]{geometry}
\usepackage{booktabs,tabularx,array}
\usepackage{enumitem}
\newcolumntype{L}{>{\raggedright\arraybackslash}X}
\setlist{nosep, leftmargin=1.2em}
\setlength{\parskip}{3pt}
\setlength{\parindent}{0pt}
\pagestyle{empty}

\begin{document}

{\large\textbf{Anexo técnico: uso de inteligencia artificial}}\\
Proyecto 2 --- Visualización I · Arnold Santiago Torres Anzola, Juan David Cruz Romero, Daniel Alejandro Cipagauta Palacios

\section*{1. Herramienta y propósito}
\textbf{Herramienta:} Claude (Anthropic, modelo Claude Opus 5.5) a través de la extensión Claude Code para VS Code.
\textbf{Propósito:} apoyar la reformulación de la audiencia y la pregunta según la retroalimentación de la docente;
inspeccionar el archivo oficial; generar y documentar el código del Notebook; proponer la paleta de colores;
diseñar el informe HTML de Plotly; redactar borradores del documento PDF y de la presentación de ensayo; y subir
los archivos al repositorio. El equipo definió el enfoque, aprobó cada decisión y es responsable del resultado.

\section*{2. Prompts empleados}
\small
\begin{tabularx}{\textwidth}{@{}p{7.3cm}L@{}}
\toprule
\textbf{Prompt (resumido)} & \textbf{Para qué se usó} \\
\midrule
``Ayúdame a corregir este archivo [propuesta] con estas especificaciones [retroalimentación de la docente] y dime cómo se podría hacer'' & Delimitar la audiencia, reformular la pregunta desde el fenómeno y eliminar el análisis por edades, que el archivo no permite. \\
``Mira, este es el archivo de los datos [WHR24\_Data\_Figure\_2.1.xls]'' & Verificar columnas, faltantes y que los componentes suman la evaluación de vida. \\
``Ayúdame a generar la pregunta problema'' · ``Explícame la opción 1'' & Comparar tres formulaciones y elegir la comparación Q4 frente a Q1. \\
``¿Qué gráficas utilizarías?'' & Seleccionar 4 gráficos de Seaborn y 2 de Plotly con propósito analítico y \textbf{proponer la paleta} (ver sección 3). \\
``Dale, hazlo'' & Generar y ejecutar el Notebook completo con interpretaciones. \\
``Explícame qué hiciste en el código'' & Comprender cada sección para la sustentación. \\
``Ve al HTML, metámosle mejores cosas, está muy plano'' & Convertir la exportación de Plotly en un informe con contexto y lecturas. \\
``En este te falta los principales hallazgos'' · ``Dame una presentación de esto'' & Completar el documento PDF y preparar diapositivas de ensayo con guion. \\
\bottomrule
\end{tabularx}
\normalsize

\section*{3. Paleta de colores}
No se escribió un prompt dedicado a la paleta: la IA la propuso al responder ``¿Qué gráficas utilizarías?'' y el
equipo la aprobó con ``Dale, hazlo''. Criterios propuestos y aceptados:
\begin{itemize}
\item \textbf{Cuartiles:} un solo tono azul de claro (Q1) a oscuro (Q4), porque los grupos tienen orden.
\item \textbf{Dimensiones:} una paleta categórica de referencia, validada para daltonismo en pares adyacentes, en orden fijo; cada componente conserva su color en todos los gráficos. \textit{Distopía + residual} en gris, por ser lo que el modelo no explica.
\item \textbf{Correlaciones:} escala divergente rojo--gris--azul con gris en el cero.
\end{itemize}
Limitación: el validador automático de la paleta no pudo ejecutarse en el equipo (requería Node.js), por lo que
se usó la paleta ya validada sin modificar sus colores; los colores de bajo contraste se acompañan de leyendas y etiquetas.

\section*{4. Revisión crítica}
\begin{itemize}
\item \textbf{Fuente de datos:} se detectó que el primer archivo usado era la versión 2026 y se reemplazó por el
  archivo oficial WHR24, como exige la guía.
\item \textbf{Errores corregidos:} el diccionario de la propuesta original describía valores de indicadores que el
  archivo no contiene (solo trae contribuciones); un título de gráfico tenía una cifra escrita a mano que no coincidía con el
  cálculo (0,91 en lugar de 0,90); una interpretación exageraba un dato del Líbano; Venezuela (PIB = 0) distorsionaba la
  recta usada para identificar atípicos y se excluyó de ese ajuste.
\item \textbf{Cautela inferencial:} se revisaron títulos e interpretaciones para eliminar lenguaje causal
  (por ejemplo, ``la dimensión socioinstitucional los separa'' se cambió por una descripción).
\item \textbf{Verificación:} todas las cifras del informe salen del Notebook, que se ejecuta de principio a fin;
  los integrantes revisaron el código y las interpretaciones para poder explicarlos en la sustentación.
\item \textbf{Reflexión:} la IA aceleró la programación y el diseño, pero también cometió errores que solo
  aparecieron al contrastar los resultados con los datos. Su uso exige verificar cada afirmación contra la evidencia.
\end{itemize}

\end{document}
